% Part: incompleteness
% Chapter: arithmetization-syntax
% Section: introduction

\documentclass[../../../include/open-logic-section]{subfiles}

\begin{document}

\olfileid{inc}{art}{int}
\olsection{परिचय}

संगणनक्षमता आणि तर्कशास्त्र यांचा संबंध जोडण्यासाठी तर्कशास्त्रातील
वस्तू (चिन्हे, पदे, !!{formula}s, !!{derivation}s), त्यांवरील क्रिया,
आणि त्यांचे गुणधर्म व संबंध यांविषयी संगणकीय हाताळणीस अनुकूल रीतीने
बोलता आले पाहिजे. चिन्हे, चिन्हांचे अनुक्रम आणि त्यांपासून तयार
होणाऱ्या इतर वस्तूंवरील संगणनक्षम फलने व संबंध विचारात घेऊन हे थेट
करता येते. तार्किक विन्यासमीमांसेतील सर्व वस्तू सांत असतात आणि
चिन्हांच्या !!a{enumerable} संचांपासून बनतात; म्हणून काही संगणन-प्रतिमानांत
अशी थेट हाताळणी शक्य आहे. पण पुनरावर्ती फलनांसारखी इतर संगणन-प्रतिमाने
संख्या आणि त्यांवरील संबंध व फलने यांपुरती मर्यादित असतात. शिवाय,
अखेरीस काही उपपत्त्यांच्या आतही विन्यासमीमांसा हाताळायची आहे;
विशेषतः अंकगणिताच्या भाषेत मांडलेल्या उपपत्त्यांमध्ये. अशा वेळी
विन्यासमीमांसेचे \emph{अंकगणितीकरण} करावे लागते: विन्यासात्मक वस्तू,
त्यांवरील क्रिया आणि त्यांचे संबंध यांचे अनुक्रमे संख्या, अंकगणितीय
फलने आणि अंकगणितीय संबंध यांच्या रूपात निरूपण करावे लागते. लाइब्निझपर्यंत
मागे जाणारी यामागील कल्पना म्हणजे विन्यासात्मक वस्तूंना संख्या नेमून देणे.

चिन्हांना त्यांचे ``संकेतांक'' म्हणून संख्या देणे तुलनेने सोपे आहे.
काही चिन्हांमुळे थोडी अडचण येते; उदाहरणार्थ, !!{variable}s अनंत असतात
आणि प्रत्येक स्थानसंख्या~$n$ साठी !!{function}s देखील अनंत असतात.
तरीही चिन्हांना पद्धतशीरपणे संख्या देता येतात, जेणेकरून, उदाहरणार्थ,
$\Obj v_2$ आणि $\Obj v_3$ यांना वेगवेगळे संकेतांक मिळतील. चिन्हांचे
अनुक्रम (उदा. पदे आणि !!{formula}s) सांकेतिक करणे अधिक अवघड आहे.
पण संख्यांचे अनुक्रम निव्वळ अंकगणितीय रीतीने हाताळता आले (उदा.
अविभाज्य संख्यांच्या घातांनी अनुक्रम सांकेतिक करून), तर एकेका चिन्हाचे
सांकेतीकरण चिन्हांच्या अनुक्रमांपर्यंत वाढवता येते; त्यानंतर
!!{formula}s चे अनुक्रम किंवा इतर मांडण्या, जसे !!{derivation}s,
यांपर्यंतही ते वाढवता येते. या विस्तारित सांकेतीकरणाला ``गोडेल
क्रमांकन'' म्हणतात. प्रत्येक पदाला, !!{formula} ला आणि
!!{derivation} ला गोडेल संख्या नेमली जाते.

चिन्हांचे अनुक्रम त्यांच्या संकेतांकांच्या अनुक्रमांनी सांकेतिक करून,
आणि संगणनक्षम फलनांनी हाताळता येईल अशी अनुक्रम-सांकेतीकरणाची पद्धत
निवडून, गोडेल संख्याही संगणनक्षम फलनांनी हाताळता येतात. प्रत्यक्षात
येथे लागणारी सर्व फलने आदिम पुनरावर्ती असतील. उदाहरणार्थ, अनुक्रमाच्या
संकेतांकावरून त्याची लांबी आणि त्यातील $i$-वा घटक काढण्याची दोन्ही
फलने आदिम पुनरावर्ती आहेत. अनुक्रम सांकेतिक करणारी संख्या, उदाहरणार्थ,
!!a{formula}~$!A$ ची गोडेल संख्या असेल, तर !!a{formula} ची लांबी
आणि !!a{formula} मधील $i$-व्या चिन्हाचा संकेतांक हेही~$!A$ च्या
गोडेल संख्येवरून काढता येतात, हे लगेच दिसते. एखादी संख्या योग्य
रीतीने रचलेल्या पदाची किंवा योग्य !!{derivation} ची गोडेल संख्या
आहे हा गुणधर्म आदिम पुनरावर्ती आहे, हे सिद्ध करणे थोडे अधिक अवघड आहे.
तरीही ते शक्य आहे, कारण आपल्याला अभिप्रेत अनुक्रमांची (पदे,
!!{formula}s, !!{derivation}s) व्याख्या विगमनाधारित रीतीने केलेली आहे.

उदाहरण म्हणून आदेशनाची क्रिया विचारात घेऊ. $!A$ हे सूत्र, $x$ हा
चर आणि $t$ हे पद असेल, तर $!A$ मधील~$x$ च्या प्रत्येक मुक्त
आस्थितीऐवजी~$t$ ठेवल्यावर $\Subst{!A}{t}{x}$ मिळते. आता $!A$,
$x$, $t$ यांना अनुक्रमे $k$, $l$, $m$ या गोडेल संख्या दिल्या आहेत
असे समजू. तीच पद्धत $\Subst{!A}{t}{x}$ लाही एखादी गोडेल संख्या,
समजा~$n$, देते. $k$, $l$, $m$ यांना~$n$ शी जोडणारे हे फलन म्हणजे
आदेशन-क्रियेचे अंकगणितीय समरूप आहे. आदेशन-क्रिया $!A$, $x$, $t$
यांपासून $\Subst{!A}{t}{x}$ देते, तर अंकगणितीकृत आदेशन-फलन
$k$, $l$, $m$ या गोडेल संख्यांपासून गोडेल संख्या~$n$ देते.
हे फलन आदिम पुनरावर्ती आहे, हे आपण पाहू.

विन्यासमीमांसेचे अंकगणितीकरण केवळ अमूर्त रुचीचे नाही. अंकगणिताची
भाषा यांसारख्या भाषांत भाषांविषयी ``बोलण्याची'' साधने आधीपासून
नसूनही व्यंजकांचे गुंतागुंतीचे गुणधर्म त्यांत आकारिक रीतीने व्यक्त
करता येतात, ही कल्पना सुरुवातीला सहज सुचण्यासारखी नव्हती. मग
अशा भाषेतील, उदाहरणार्थ पेआनो अंकगणितासारखी, उपपत्ती स्वतःच्या
भाषेविषयी काय \emph{सिद्ध} करू शकते, हा प्रश्न विचारणे स्वाभाविक
ठरते (उदा. !!{sentence}s सिद्ध करता येतात का किंवा ती सत्य आहेत
का). यातून गोडेलची (असिद्धतायोग्यतेविषयी) आणि टार्स्कीची
(सत्याच्या अव्याख्येयतेविषयी) मर्यादा दाखवणारी प्रसिद्ध प्रमेये
समोर येतात. पण संगणनक्षमता उपपत्तीतील काही महत्त्वाचे निष्कर्ष
सिद्ध करण्यासाठीही विन्यासमीमांसेचे अंकगणितीकरण महत्त्वाचे आहे;
उदा. उपपत्त्यांची संगणकीय शक्ती किंवा संगणनक्षमतेच्या विविध
प्रतिमानांतील संबंध यांविषयीचे निष्कर्ष. इतर वस्तू आणि गुणधर्मांचे
अंकगणितीकरण कसे करावे याचा नमुनाही विन्यासमीमांसेचे अंकगणितीकरण
देते. उदाहरणार्थ, (ट्यूरिंग यंत्रांच्या) स्थितिवर्णनांचे आणि
संगणनांचेही अशाच रीतीने अंकगणितीकरण करता येते. त्यामुळे एका
प्रतिमानातील संगणनांचे (उदा. ट्यूरिंग यंत्रांतील) दुसऱ्या प्रतिमानात
(उदा. पुनरावर्ती फलनांत) अनुकरण करता येते.

\end{document}
